% Source: https://artofproblemsolving.com/community/c3h1663330_clock_problem
% Stable identity: aops-topic-1663330-q01
Measure positions on the dial in ``minute marks'' (there are $60$ around the dial). At $4{:}00$ the minute hand is at mark $0$ and the hour hand is at mark $20$.

In $x$ minutes the minute hand moves $x$ marks. The hour hand moves from one hour number to the next, that is $5$ marks, every $60$ minutes, so in $x$ minutes it moves $5\cdot\dfrac{x}{60}$ marks. This is where the $5$ comes from. The hands coincide when the minute hand reaches the hour hand:
\[
x=20+5\cdot\dfrac{x}{60}=20+\dfrac{x}{12}\qquad\Longrightarrow\qquad \dfrac{11x}{12}=20\qquad\Longrightarrow\qquad x=\dfrac{240}{11}.
\]
So the hands coincide after
\[
\boxed{\dfrac{240}{11}=21\dfrac{9}{11}\text{ minutes}},
\]
at about $4{:}21{:}49$. (Equivalently, the minute hand gains $60-5=55$ marks per hour, that is $\dfrac{11}{12}$ mark per minute, on the hour hand, and it must gain $20$ marks.)
